% Options for packages loaded elsewhere
\PassOptionsToPackage{unicode}{hyperref}
\PassOptionsToPackage{hyphens}{url}
\PassOptionsToPackage{dvipsnames,svgnames,x11names}{xcolor}
\documentclass[
  letterpaper,
]{article}
\usepackage{xcolor}
\usepackage[margin=1in]{geometry}
\usepackage{amsmath,amssymb}
\setcounter{secnumdepth}{-\maxdimen} % remove section numbering
\usepackage{iftex}
\ifPDFTeX
  \usepackage[T1]{fontenc}
  \usepackage[utf8]{inputenc}
  \usepackage{textcomp} % provide euro and other symbols
\else % if luatex or xetex
  \usepackage{unicode-math} % this also loads fontspec
  \defaultfontfeatures{Scale=MatchLowercase}
  \defaultfontfeatures[\rmfamily]{Ligatures=TeX,Scale=1}
\fi
\usepackage{lmodern}
\ifPDFTeX\else
  % xetex/luatex font selection
\fi
% Use upquote if available, for straight quotes in verbatim environments
\IfFileExists{upquote.sty}{\usepackage{upquote}}{}
\IfFileExists{microtype.sty}{% use microtype if available
  \usepackage[]{microtype}
  \UseMicrotypeSet[protrusion]{basicmath} % disable protrusion for tt fonts
}{}
\makeatletter
\@ifundefined{KOMAClassName}{% if non-KOMA class
  \IfFileExists{parskip.sty}{%
    \usepackage{parskip}
  }{% else
    \setlength{\parindent}{0pt}
    \setlength{\parskip}{6pt plus 2pt minus 1pt}}
}{% if KOMA class
  \KOMAoptions{parskip=half}}
\makeatother
\usepackage{color}
\usepackage{fancyvrb}
\newcommand{\VerbBar}{|}
\newcommand{\VERB}{\Verb[commandchars=\\\{\}]}
\DefineVerbatimEnvironment{Highlighting}{Verbatim}{commandchars=\\\{\}}
% Add ',fontsize=\small' for more characters per line
\newenvironment{Shaded}{}{}
\newcommand{\AlertTok}[1]{\textcolor[rgb]{1.00,0.00,0.00}{\textbf{#1}}}
\newcommand{\AnnotationTok}[1]{\textcolor[rgb]{0.38,0.63,0.69}{\textbf{\textit{#1}}}}
\newcommand{\AttributeTok}[1]{\textcolor[rgb]{0.49,0.56,0.16}{#1}}
\newcommand{\BaseNTok}[1]{\textcolor[rgb]{0.25,0.63,0.44}{#1}}
\newcommand{\BuiltInTok}[1]{\textcolor[rgb]{0.00,0.50,0.00}{#1}}
\newcommand{\CharTok}[1]{\textcolor[rgb]{0.25,0.44,0.63}{#1}}
\newcommand{\CommentTok}[1]{\textcolor[rgb]{0.38,0.63,0.69}{\textit{#1}}}
\newcommand{\CommentVarTok}[1]{\textcolor[rgb]{0.38,0.63,0.69}{\textbf{\textit{#1}}}}
\newcommand{\ConstantTok}[1]{\textcolor[rgb]{0.53,0.00,0.00}{#1}}
\newcommand{\ControlFlowTok}[1]{\textcolor[rgb]{0.00,0.44,0.13}{\textbf{#1}}}
\newcommand{\DataTypeTok}[1]{\textcolor[rgb]{0.56,0.13,0.00}{#1}}
\newcommand{\DecValTok}[1]{\textcolor[rgb]{0.25,0.63,0.44}{#1}}
\newcommand{\DocumentationTok}[1]{\textcolor[rgb]{0.73,0.13,0.13}{\textit{#1}}}
\newcommand{\ErrorTok}[1]{\textcolor[rgb]{1.00,0.00,0.00}{\textbf{#1}}}
\newcommand{\ExtensionTok}[1]{#1}
\newcommand{\FloatTok}[1]{\textcolor[rgb]{0.25,0.63,0.44}{#1}}
\newcommand{\FunctionTok}[1]{\textcolor[rgb]{0.02,0.16,0.49}{#1}}
\newcommand{\ImportTok}[1]{\textcolor[rgb]{0.00,0.50,0.00}{\textbf{#1}}}
\newcommand{\InformationTok}[1]{\textcolor[rgb]{0.38,0.63,0.69}{\textbf{\textit{#1}}}}
\newcommand{\KeywordTok}[1]{\textcolor[rgb]{0.00,0.44,0.13}{\textbf{#1}}}
\newcommand{\NormalTok}[1]{#1}
\newcommand{\OperatorTok}[1]{\textcolor[rgb]{0.40,0.40,0.40}{#1}}
\newcommand{\OtherTok}[1]{\textcolor[rgb]{0.00,0.44,0.13}{#1}}
\newcommand{\PreprocessorTok}[1]{\textcolor[rgb]{0.74,0.48,0.00}{#1}}
\newcommand{\RegionMarkerTok}[1]{#1}
\newcommand{\SpecialCharTok}[1]{\textcolor[rgb]{0.25,0.44,0.63}{#1}}
\newcommand{\SpecialStringTok}[1]{\textcolor[rgb]{0.73,0.40,0.53}{#1}}
\newcommand{\StringTok}[1]{\textcolor[rgb]{0.25,0.44,0.63}{#1}}
\newcommand{\VariableTok}[1]{\textcolor[rgb]{0.10,0.09,0.49}{#1}}
\newcommand{\VerbatimStringTok}[1]{\textcolor[rgb]{0.25,0.44,0.63}{#1}}
\newcommand{\WarningTok}[1]{\textcolor[rgb]{0.38,0.63,0.69}{\textbf{\textit{#1}}}}
\usepackage{longtable,booktabs,array}
\newcounter{none} % for unnumbered tables
\usepackage{calc} % for calculating minipage widths
% Correct order of tables after \paragraph or \subparagraph
\usepackage{etoolbox}
\makeatletter
\patchcmd\longtable{\par}{\if@noskipsec\mbox{}\fi\par}{}{}
\makeatother
% Allow footnotes in longtable head/foot
\IfFileExists{footnotehyper.sty}{\usepackage{footnotehyper}}{\usepackage{footnote}}
\makesavenoteenv{longtable}
\setlength{\emergencystretch}{3em} % prevent overfull lines
\providecommand{\tightlist}{%
  \setlength{\itemsep}{0pt}\setlength{\parskip}{0pt}}
\usepackage{bookmark}
\IfFileExists{xurl.sty}{\usepackage{xurl}}{} % add URL line breaks if available
\urlstyle{same}
\hypersetup{
  pdftitle={Parallel Stateful Agent Search},
  colorlinks=true,
  linkcolor={blue},
  filecolor={Maroon},
  citecolor={Blue},
  urlcolor={blue},
  pdfcreator={LaTeX via pandoc}}

\title{Parallel Stateful Agent Search}
\usepackage{etoolbox}
\makeatletter
\providecommand{\subtitle}[1]{% add subtitle to \maketitle
  \apptocmd{\@title}{\par {\large #1 \par}}{}{}
}
\makeatother
\subtitle{A Systems and Research Assessment}
\author{}
\date{}

\begin{document}
\maketitle

{
\setcounter{tocdepth}{3}
\tableofcontents
}
\section{Bottom line}\label{bottom-line}

\textbf{Yes, conditionally.} Parallel agent forking can reduce \textbf{time-to-verified-solution}, but it does not reduce the causal depth of the correct solution. It helps by overlapping work that a sequential agent would otherwise spend on:

\begin{itemize}
\tightlist
\item
  wrong hypotheses and restarts,
\item
  uncertain high-impact choices,
\item
  independent evidence gathering or tool calls,
\item
  heavy-tailed failures,
\item
  parallelizable subtasks.
\end{itemize}

It is unlikely to help when the task has one obvious path, is mostly sequential, branches are highly correlated, tools are bottlenecked, or intermediate progress cannot be evaluated reliably.

The strongest version of the idea is not ``clone the same chat \(N\) times.'' It is:

\begin{quote}
\textbf{Asynchronous, cost-aware search over checkpointed agent execution states, using deliberate diversity, predictive evaluation, adaptive horizons, and transactional side-effect isolation.}
\end{quote}

The search structure will usually be an \textbf{AND/OR DAG with stochastic transitions}, not a clean tree.

\begin{center}\rule{0.5\linewidth}{0.5pt}\end{center}

\section{1. What exactly is being forked?}\label{what-exactly-is-being-forked}

A stateful-agent node is much larger than a text prefix:

\[
s = (\text{task}, \text{transcript}, \text{controller/memory},
\text{workspace}, \text{tool/world state}, \text{pending effects},
\text{model/config}, \text{budget}, \text{lineage})
\]

A transcript alone is insufficient: two agents with identical text but different files, browser sessions, database versions, pending calls, or credentials are not in the same state.

A practical architecture is:

\begin{Shaded}
\begin{Highlighting}[]
\NormalTok{Task/specification}
\NormalTok{       |}
\NormalTok{Immutable root checkpoint}
\NormalTok{       |}
\NormalTok{Proposal generator: create K \textgreater{} N structured alternatives}
\NormalTok{       |}
\NormalTok{Diversity + value portfolio selector}
\NormalTok{       |}
\NormalTok{Copy{-}on{-}write sandbox workers, up to N concurrently}
\NormalTok{       |}
\NormalTok{Roll out each branch to adaptive horizon X}
\NormalTok{       |}
\NormalTok{Hard tests + value model + uncertainty + deduplication}
\NormalTok{       |}
\NormalTok{Async frontier scheduler}
\NormalTok{   /       |       \textbackslash{}}
\NormalTok{prune   continue   fork again}
\NormalTok{       |}
\NormalTok{Select / typed merge / synthesis}
\NormalTok{       |}
\NormalTok{Clean replay, independent verification, single commit}
\end{Highlighting}
\end{Shaded}

Each proposal should describe:

\begin{Shaded}
\begin{Highlighting}[]
\NormalTok{assumptions}
\NormalTok{strategy or hypothesis}
\NormalTok{planned actions}
\NormalTok{expected milestone at horizon X}
\NormalTok{falsifying observation}
\NormalTok{estimated cost}
\NormalTok{side{-}effect class}
\end{Highlighting}
\end{Shaded}

This makes branches comparable before spending a full rollout on them.

\begin{center}\rule{0.5\linewidth}{0.5pt}\end{center}

\section{2. Does parallel forking reduce wall-clock time?}\label{does-parallel-forking-reduce-wall-clock-time}

\subsection{The important distinction}\label{the-important-distinction}

If one sequential run always takes \(\tau\), and \(N\) replicas also each take \(\tau\), running all \(N\) does \textbf{not} make a result appear before \(\tau\). It raises the probability that a good result exists at time \(\tau\).

Actual latency reduction occurs when parallelism overlaps:

\begin{enumerate}
\def\labelenumi{\arabic{enumi}.}
\tightlist
\item
  attempts that would otherwise be tried sequentially;
\item
  variable-time trajectories, allowing stop-on-first-success;
\item
  backtracking and debugging;
\item
  independent tool latency;
\item
  genuinely separable subtasks.
\end{enumerate}

\subsection{Work--span view}\label{workspan-view}

Let:

\begin{itemize}
\tightlist
\item
  \(W(N,X)\): total work across all branches, including evaluation and canceled work;
\item
  \(P\): actual concurrent execution capacity;
\item
  \(D(N,X)\): longest causal chain of the selected solution;
\item
  \(H(N,X)\): snapshot, queue, coordination, evaluation, merge, and commit overhead.
\end{itemize}

Any implementation is bounded by approximately:

\[
T_{\text{parallel}} \gtrsim
\max\left(\frac{W(N,X)}{P}, D(N,X)\right)+H(N,X)
\]

Therefore, a necessary feasibility condition is:

\[
\max\left(\frac{W}{P},D\right)+H<T_{\text{sequential}}
\]

Parallelism cannot reduce \(D\). It can only remove sequential search and wrong turns around that critical path.

\subsection{Simple independent-attempt model}\label{simple-independent-attempt-model}

Suppose one branch reaches an acceptable result by horizon \(X\) with probability \(p_X\). Under independence:

\[
P_{\text{any}}(N,X)=1-(1-p_X)^N
\]

If the coordinator successfully recognizes or selects a good branch with probability \(\eta(N,X)\), then approximately:

\[
P_{\text{delivered}}(N,X)
=
\eta(N,X)\left[1-(1-p_X)^N\right]
\]

This distinction is critical: \textbf{generating a correct branch is not the same as returning it}.

Let:

\begin{itemize}
\tightlist
\item
  \(R_N(X)\): wall time of one parallel batch, including overhead;
\item
  \(W_N(X)\): compute cost of one batch.
\end{itemize}

If batches are repeated until a verified success:

\[
\mathbb E[T_{\text{parallel}}]
=
\frac{R_N(X)}{P_{\text{delivered}}(N,X)}
\]

\[
\mathbb E[W_{\text{parallel}}]
=
\frac{W_N(X)}{P_{\text{delivered}}(N,X)}
\]

For one sequential restart policy:

\[
\mathbb E[T_{\text{sequential}}]
=
\frac{R_1(X)}{p_X}
\]

Thus parallelism wins on expected latency when:

\[
\boxed{
\frac{P_{\text{delivered}}(N,X)}{R_N(X)}
>
\frac{p_X}{R_1(X)}
}
\]

That is: \textbf{verified successes per wall-clock second must increase}.

An economic constraint can be added:

\[
\boxed{
\frac{W_N(X)}{P_{\text{delivered}}(N,X)}
\le
\kappa\frac{W_1(X)}{p_X}
}
\]

where \(\kappa\) is the acceptable compute multiplier.

\subsubsection{Toy example}\label{toy-example}

If \(p_X=0.1\), \(N=4\), branches are independent, and batch overhead is \(0.1\tau\):

\[
P_{\text{any}}=1-0.9^4=0.3439
\]

Expected sequential retry time is \(10\tau\). Expected parallel batch time is:

\[
\frac{1.1\tau}{0.3439}\approx3.2\tau
\]

This is about a \(3.1\times\) latency improvement for roughly \(16\%\) more expected branch work per success. This is a favorable rare-success regime.

When \(p_X\) is already high, additional replicas usually buy little latency and much more compute.

\subsection{Correlation changes everything}\label{correlation-changes-everything}

The independence formula is optimistic. Same-model branches share:

\begin{itemize}
\tightlist
\item
  training-data blind spots,
\item
  prompt anchoring,
\item
  tool biases,
\item
  evaluator biases,
\item
  preferred implementation patterns.
\end{itemize}

A common-mode failure probability creates a success ceiling no matter how large \(N\) becomes. Therefore the useful quantity is not nominal \(N\), but the empirical \textbf{conditional marginal success} of the next branch.

\subsection{Scaling behavior}\label{scaling-behavior}

{\def\LTcaptype{none} % do not increment counter
\begin{longtable}[]{@{}
  >{\raggedright\arraybackslash}p{(\linewidth - 4\tabcolsep) * \real{0.3333}}
  >{\raggedright\arraybackslash}p{(\linewidth - 4\tabcolsep) * \real{0.3333}}
  >{\raggedright\arraybackslash}p{(\linewidth - 4\tabcolsep) * \real{0.3333}}@{}}
\toprule\noalign{}
\begin{minipage}[b]{\linewidth}\raggedright
Quantity
\end{minipage} & \begin{minipage}[b]{\linewidth}\raggedright
As \(N\) increases
\end{minipage} & \begin{minipage}[b]{\linewidth}\raggedright
As \(X\) increases
\end{minipage} \\
\midrule\noalign{}
\endhead
\bottomrule\noalign{}
\endlastfoot
Success coverage & Increases, then saturates & Usually increases \\
Total work & Roughly \(O(NX)\) per round & Increases per losing branch \\
Wall latency & Falls toward critical path, then flattens or rises & Longer rollouts, fewer coordination rounds \\
Evaluator difficulty & More multiple-comparison errors & Usually more informative states \\
Coordination & \(O(N)\) scoring; naive pairwise \(O(N^2)\) & Fewer evaluations \\
Diversity benefit & Diminishing with correlation & Branches may diverge more \\
Pruning risk & More alternatives to rank & Delayed-payoff paths become more visible \\
\end{longtable}
}

Waiting for every branch produces a straggler-sensitive \(\max_i L_i\). An asynchronous stop-on-first-verified-success or top-\(k\) quorum is generally better.

\begin{center}\rule{0.5\linewidth}{0.5pt}\end{center}

\section{3. What should agents diversify over?}\label{what-should-agents-diversify-over}

Diversify over \textbf{consequential decisions}, not writing style.

High-value axes include:

\begin{enumerate}
\def\labelenumi{\arabic{enumi}.}
\tightlist
\item
  \textbf{Task interpretation and assumptions}

  \begin{itemize}
  \tightlist
  \item
    competing readings of ambiguous requirements;
  \item
    different constraint priorities.
  \end{itemize}
\item
  \textbf{Decomposition}

  \begin{itemize}
  \tightlist
  \item
    alternative subproblem boundaries;
  \item
    different subtask ordering;
  \item
    OR alternatives versus AND-composable subtasks.
  \end{itemize}
\item
  \textbf{Causal hypothesis}

  \begin{itemize}
  \tightlist
  \item
    suspected root causes;
  \item
    proof lemmas;
  \item
    failure mechanisms.
  \end{itemize}
\item
  \textbf{Strategy}

  \begin{itemize}
  \tightlist
  \item
    proof-first versus experiment-first;
  \item
    minimal patch versus architectural correction;
  \item
    retrieval-first versus direct implementation.
  \end{itemize}
\item
  \textbf{Tool and evidence acquisition}

  \begin{itemize}
  \tightlist
  \item
    different queries, sources, tests, inspections, simulators, or environments.
  \end{itemize}
\item
  \textbf{Implementation}

  \begin{itemize}
  \tightlist
  \item
    algorithms, data structures, APIs, libraries, or patch designs.
  \end{itemize}
\item
  \textbf{Predicted horizon state}

  \begin{itemize}
  \tightlist
  \item
    expected artifact;
  \item
    expected test result;
  \item
    milestone and falsifier.
  \end{itemize}
\item
  \textbf{Model or controller heterogeneity}

  \begin{itemize}
  \tightlist
  \item
    different models, prompts, temperatures, tool policies, or context views.
  \end{itemize}
\end{enumerate}

Random seeds alone often produce lexical rather than behavioral diversity.

\subsection{Portfolio selection}\label{portfolio-selection}

Generate \(K\gg N\) cheap branch proposals, cluster them, then choose a quality-diversity portfolio such as:

\[
A^*=
\arg\max_{|A|=N}
\left[
\sum_{i\in A}\mu_i
-
\lambda\sum_{i<j}\operatorname{sim}(i,j)
\right]
\]

Here \(\mu_i\) is predicted value and similarity should use:

\begin{itemize}
\tightlist
\item
  assumptions,
\item
  plan graphs,
\item
  proposed tool calls,
\item
  expected artifacts,
\item
  code or action signatures,
\end{itemize}

not merely text embeddings.

A reasonable portfolio includes:

\begin{itemize}
\tightlist
\item
  several high-value exploit branches;
\item
  one branch per materially different strategy cluster;
\item
  a small wildcard/exploration quota.
\end{itemize}

Keep branches private through their first horizon. Immediate shared memory often causes herding.

\begin{center}\rule{0.5\linewidth}{0.5pt}\end{center}

\section{\texorpdfstring{4. How should horizon \(X\) be defined?}{4. How should horizon X be defined?}}\label{how-should-horizon-x-be-defined}

A universal token horizon is usually a poor choice. Use a \textbf{semantic stopping time plus resource ceilings}.

Represent the cap as:

\[
X=(\text{actions},\text{tool calls},\text{tokens},
\text{worker-seconds},\text{dollars},\text{effect budget})
\]

A rollout should normally stop at the first of:

\begin{itemize}
\tightlist
\item
  a new informative tool observation;
\item
  a test, compile, or simulation result;
\item
  completion of a plan milestone;
\item
  hard failure or loop detection;
\item
  an irreversible side-effect boundary;
\item
  a resource ceiling;
\item
  sufficiently strong evidence for pruning or promotion.
\end{itemize}

For coding, one useful macro-action might be:

\begin{quote}
inspect or hypothesize → edit → run targeted test → assimilate result → checkpoint.
\end{quote}

\subsection{Short versus long horizons}\label{short-versus-long-horizons}

\textbf{Too short:}

\begin{itemize}
\tightlist
\item
  evaluations are noisy;
\item
  agents cannot establish coherent progress;
\item
  coordination dominates;
\item
  good delayed-payoff paths are repeatedly pruned.
\end{itemize}

\textbf{Too long:}

\begin{itemize}
\tightlist
\item
  compute is wasted on bad branches;
\item
  mistakes compound;
\item
  correction comes late;
\item
  fewer alternatives are tested.
\end{itemize}

If a solution path has effective depth \(D\), and pruning occurs about \(D/X\) times, then even good per-stage retention can compound badly. If the probability of retaining the correct continuation at each stage is \(r\):

\[
P_{\text{survive}}\approx r^{D/X}
\]

Thus many short horizons can be worse than a few longer ones.

A good default is asynchronous successive halving:

\begin{Shaded}
\begin{Highlighting}[]
\NormalTok{many branches receive a small pilot}
\NormalTok{promising or uncertain branches receive 2X}
\NormalTok{survivors receive 4X}
\NormalTok{finalists receive enough budget to finish}
\end{Highlighting}
\end{Shaded}

Branches should be compared at similar budget rungs. Horizon should become longer when value estimates are uncertain or rewards are delayed.

\begin{center}\rule{0.5\linewidth}{0.5pt}\end{center}

\section{5. Predicting which branch is worth continuing}\label{predicting-which-branch-is-worth-continuing}

There is no general cheap predictor: on open-ended tasks, determining whether a branch will succeed may be nearly as difficult as finishing it. The practical solution is a cascade.

\subsection{Evaluation cascade}\label{evaluation-cascade}

\begin{enumerate}
\def\labelenumi{\arabic{enumi}.}
\tightlist
\item
  \textbf{Hard gates}

  \begin{itemize}
  \tightlist
  \item
    syntax, schema, constraints;
  \item
    compile and tests;
  \item
    safety policy;
  \item
    duplicated state or repeated failure;
  \item
    forbidden side effects.
  \end{itemize}
\item
  \textbf{Cheap progress features}

  \begin{itemize}
  \tightlist
  \item
    verified requirements covered;
  \item
    newly passing tests;
  \item
    unresolved blockers;
  \item
    patch complexity;
  \item
    quality of evidence;
  \item
    tool success history.
  \end{itemize}
\item
  \textbf{Learned process/value model}

  \begin{itemize}
  \tightlist
  \item
    predict:
    \[
    V(s,b)=P(\text{terminal verifier passes by deadline}\mid s,b)
    \]
  \item
    also predict remaining cost and uncertainty.
  \end{itemize}
\item
  \textbf{Pairwise or tournament evaluation}

  \begin{itemize}
  \tightlist
  \item
    useful for close finalists;
  \item
    avoid all-pairs \(O(N^2)\) comparisons.
  \end{itemize}
\item
  \textbf{Short Monte Carlo rollouts}

  \begin{itemize}
  \tightlist
  \item
    execute several cheap continuations from a promising state;
  \item
    use a smaller model or reduced tool environment where possible.
  \end{itemize}
\item
  \textbf{Full verifier}

  \begin{itemize}
  \tightlist
  \item
    expensive tests, simulation, theorem checker, or blinded reviewer.
  \end{itemize}
\end{enumerate}

\subsection{Value of computation}\label{value-of-computation}

Let \(v^*\) be the incumbent value. Continue branch \(i\) for another quantum \(\delta\) when:

\[
\operatorname{VOC}_i(\delta)=
\mathbb E[\max(v^*,V_i^{\text{after }\delta})]-v^*
-\lambda_c\mathbb E[C_i]
-\lambda_t\mathbb E[T_i]
>0
\]

Possible schedulers include:

\begin{itemize}
\tightlist
\item
  UCB or Thompson sampling;
\item
  MCTS-style selection;
\item
  best-first search;
\item
  successive halving/Hyperband;
\item
  branch-and-bound when a genuinely valid upper bound exists.
\end{itemize}

A useful priority heuristic is:

\[
\frac{
\text{expected improvement}
+\beta\,\text{uncertainty}
+\gamma\,\text{information gain}
+\rho\,\text{novelty}
}{
\text{predicted incremental cost}
}
\]

\subsection{Evaluator risks}\label{evaluator-risks}

As \(N\) grows, the highest judge score increasingly contains ``winner's curse'' noise. More candidates can therefore reduce delivered quality even while oracle coverage rises.

Mitigations:

\begin{itemize}
\tightlist
\item
  retain uncertainty-aware exploration;
\item
  independently re-evaluate the winner;
\item
  use objective tests where possible;
\item
  use heterogeneous, blinded judges;
\item
  keep a random sample of apparently bad branches running offline.
\end{itemize}

That last point is important: if pruned branches are never continued, the system cannot measure false-negative pruning and can train itself into a self-confirming policy.

\begin{center}\rule{0.5\linewidth}{0.5pt}\end{center}

\section{6. Preventing trajectory collapse}\label{preventing-trajectory-collapse}

Use behavioral rather than textual duplicate detection.

Before execution, compare:

\begin{itemize}
\tightlist
\item
  normalized assumptions;
\item
  strategy and dependency graphs;
\item
  intended tool/action sequences;
\item
  queries and evidence sources;
\item
  expected artifacts and falsifiers.
\end{itemize}

During execution, compare:

\begin{itemize}
\tightlist
\item
  tool-call Jaccard similarity;
\item
  action-trace edit distance;
\item
  code AST or file-delta similarity;
\item
  test outcomes;
\item
  error and success correlation.
\end{itemize}

Mechanisms include:

\begin{itemize}
\tightlist
\item
  one quota per strategy cluster;
\item
  different role or model assignments;
\item
  giving already-selected proposals as negative examples;
\item
  novelty penalties;
\item
  exact transposition tables for identical states;
\item
  perturbing or canceling near-duplicates;
\item
  virtual loss or leases so workers do not simultaneously expand the same node.
\end{itemize}

Near-duplicate summaries are not enough to merge states. Exact equivalence must include workspace and tool/world versions.

\begin{center}\rule{0.5\linewidth}{0.5pt}\end{center}

\section{7. How results should be merged}\label{how-results-should-be-merged}

Different output types need different merge semantics.

{\def\LTcaptype{none} % do not increment counter
\begin{longtable}[]{@{}
  >{\raggedright\arraybackslash}p{(\linewidth - 2\tabcolsep) * \real{0.5000}}
  >{\raggedright\arraybackslash}p{(\linewidth - 2\tabcolsep) * \real{0.5000}}@{}}
\toprule\noalign{}
\begin{minipage}[b]{\linewidth}\raggedright
Situation
\end{minipage} & \begin{minipage}[b]{\linewidth}\raggedright
Preferred operation
\end{minipage} \\
\midrule\noalign{}
\endhead
\bottomrule\noalign{}
\endlastfoot
Competing global solutions & Select one coherent, end-to-end branch \\
Unique short answer from independent samples & Vote or confidence-weighted vote \\
Pairwise judgeable outputs & Tournament \\
Complementary, typed subtasks & AND-join \\
Code or documents & Three-way merge/cherry-pick, then retest \\
Facts and evidence & Provenance-preserving union \\
Mixed ideas requiring reasoning & Create a new synthesis branch \\
Identical full states & Transposition/DAG merge \\
\end{longtable}
}

A synthesis is not automatically better. Treat it as another branch that must be independently verified; otherwise it can produce an inconsistent ``Frankenstein'' result.

Shared memory should have three scopes:

\begin{enumerate}
\def\labelenumi{\arabic{enumi}.}
\tightlist
\item
  immutable root specification;
\item
  branch-local scratch state;
\item
  global ledger containing only validated facts, artifacts, failures, and provenance.
\end{enumerate}

Do not concatenate private reasoning transcripts.

\subsection{External effects}\label{external-effects}

Speculative branches should not directly send emails, purchase items, deploy code, mutate production databases, or perform other irreversible actions.

Use:

\begin{enumerate}
\def\labelenumi{\arabic{enumi}.}
\tightlist
\item
  sandboxed/dry-run execution;
\item
  side-effect intents;
\item
  idempotency keys and event logs;
\item
  prepare/revalidate;
\item
  one authoritative commit after winner selection.
\end{enumerate}

\begin{center}\rule{0.5\linewidth}{0.5pt}\end{center}

\section{8. Relation to existing methods}\label{relation-to-existing-methods}

{\def\LTcaptype{none} % do not increment counter
\begin{longtable}[]{@{}
  >{\raggedright\arraybackslash}p{(\linewidth - 4\tabcolsep) * \real{0.3333}}
  >{\raggedright\arraybackslash}p{(\linewidth - 4\tabcolsep) * \real{0.3333}}
  >{\raggedright\arraybackslash}p{(\linewidth - 4\tabcolsep) * \real{0.3333}}@{}}
\toprule\noalign{}
\begin{minipage}[b]{\linewidth}\raggedright
Existing idea
\end{minipage} & \begin{minipage}[b]{\linewidth}\raggedright
Similarity
\end{minipage} & \begin{minipage}[b]{\linewidth}\raggedright
Stateful-agent difference
\end{minipage} \\
\midrule\noalign{}
\endhead
\bottomrule\noalign{}
\endlastfoot
Speculative execution & Execute alternatives before knowing which is needed & Agent choices are heuristic and long-lived; side effects may not roll back \\
Speculative decoding & Draft futures and verify them & Token verification is cheap and exact; agent-solution verification usually is not \\
Beam search & Retain top-\(B\) prefixes & LM likelihood is not task value; agent states are large, mutable, and variable-duration \\
Tree of Thoughts / Graph of Thoughts & Search semantic reasoning units & Usually text states, not complete external execution snapshots \\
MCTS & Select, expand, rollout, back up value & Real agent environments may not be resettable, Markov, stationary, or cheap \\
Branch-and-bound / A* & Prune using incumbent and bounds & LLM judge scores are rarely admissible bounds \\
Self-consistency / best-of-\(N\) & Independent complete trajectories & No intermediate reallocation or checkpoint continuation \\
Ensembles/debate & Different perspectives followed by aggregation & Often no explicit frontier, rollback, or state search \\
Distributed portfolio solvers & Run diverse algorithms concurrently & Probably the closest top-level analogy; agent policies and world states are less structured \\
Shared-tree distributed search & Work stealing, shared frontier, transpositions & Requires concurrency control over memory, tools, and environmental effects \\
\end{longtable}
}

Empirically:

\begin{itemize}
\tightlist
\item
  \href{https://arxiv.org/abs/2203.11171}{Self-consistency} showed large gains from sampling multiple reasoning paths.
\item
  \href{https://arxiv.org/abs/2305.10601}{Tree of Thoughts} reported 74\% versus 4\% on Game of 24 for its searched versus chain-of-thought configurations.
\item
  \href{https://arxiv.org/abs/2310.04406}{LATS} applies MCTS, value estimates, reflection, and environment feedback to agents.
\item
  \href{https://arxiv.org/abs/2407.01476}{Tree Search for Language Model Agents} reported relative success gains of 39.7\% on VisualWebArena and 28\% on WebArena, but used serial best-first search and replay rather than concurrently live checkpoints.
\item
  \href{https://arxiv.org/abs/2407.21787}{Large Language Monkeys} reported SWE-bench Lite oracle coverage increasing from 15.9\% at one sample to 56\% at 250 samples. It also found selection methods plateauing without automatic verifiers.
\item
  \href{https://arxiv.org/abs/2408.03314}{Snell et al.} found that adaptive, difficulty-dependent test-time compute could be over four times as efficient as naive best-of-\(N\).
\item
  \href{https://arxiv.org/abs/2211.17192}{Speculative decoding} achieved 2--3× decoding acceleration, but relies on an exact acceptance mechanism unavailable for most agent tasks.
\end{itemize}

These results support additional test-time search, but they do \textbf{not} by themselves establish matched-compute wall-clock speedup for concurrent stateful agents.

\begin{center}\rule{0.5\linewidth}{0.5pt}\end{center}

\section{\texorpdfstring{9. When should \(T_{\text{parallel}}<T_{\text{sequential}}\)?}{9. When should T\_\{\textbackslash text\{parallel\}\}\textless T\_\{\textbackslash text\{sequential\}\}?}}\label{when-should-t_textparallelt_textsequential}

The favorable regime is:

\begin{enumerate}
\def\labelenumi{\arabic{enumi}.}
\tightlist
\item
  The success criterion is fixed and independently verifiable.
\item
  Much of sequential time is search, backtracking, waiting, or restarting rather than causal execution.
\item
  Real spare concurrency exists; the agents are not sharing a saturated GPU or rate-limited tool.
\item
  New branches add meaningful conditional diversity.
\item
  Useful progress signals appear before completion.
\item
  Evaluation is cheaper than finishing every branch.
\item
  Snapshot and cancellation overhead are small relative to \(X\).
\item
  Tools and external effects are sandboxable.
\item
  Latency has enough economic value to justify redundant compute.
\end{enumerate}

A marginal branch should be launched only when:

\[
\begin{aligned}
&V_{\text{success}}\,
\Delta P(\text{verified success before deadline})\\
&\quad+\text{value of information}\\
&>
\text{compute cost}
+\text{coordination cost}
+\text{latency externality}
+\text{safety risk}
\end{aligned}
\]

The unfavorable regime includes:

\begin{itemize}
\tightlist
\item
  strict serial dependency chains;
\item
  nearly certain first-run success;
\item
  common-mode misinterpretation;
\item
  sparse or delayed reward with weak prefix evaluation;
\item
  subjective tasks with no reliable verifier;
\item
  irreversible actions;
\item
  shared GPU/API/database bottlenecks;
\item
  tasks where a sequential agent benefits strongly from accumulating failed-attempt knowledge.
\end{itemize}

A stronger falsifiable hypothesis is:

\begin{quote}
On a preregistered distribution of sandboxable, objectively verifiable tasks with reserved concurrent capacity, an adaptive heterogeneous portfolio of at most \(N\) checkpointed branches reduces time-to-verified-solution by at least \(\delta\) relative to a tuned sequential agent, while staying within compute multiplier \(\kappa\) and maintaining non-inferior final success.
\end{quote}

\begin{center}\rule{0.5\linewidth}{0.5pt}\end{center}

\section{10. Minimal empirical experiment}\label{minimal-empirical-experiment}

\subsection{Benchmark}\label{benchmark}

Start with code repair because it offers:

\begin{itemize}
\tightlist
\item
  objective hidden tests;
\item
  meaningful multi-step tool use;
\item
  copy-on-write repository snapshots;
\item
  cheap rollback;
\item
  no real-world side effects.
\end{itemize}

Use:

\begin{itemize}
\tightlist
\item
  a stratified subset of SWE-bench Verified;
\item
  ideally an additional fresh, post-training-cutoff private set to reduce contamination.
\end{itemize}

A 40-task pilot can validate mechanics; a defensible confirmatory study should use power analysis and likely 80 or more tasks.

\subsection{Four arms}\label{four-arms}

\begin{enumerate}
\def\labelenumi{\arabic{enumi}.}
\tightlist
\item
  \textbf{S1: strong sequential}

  \begin{itemize}
  \tightlist
  \item
    one persistent agent;
  \item
    allowed planning, self-critique, backtracking, summarization, and restarts;
  \item
    receives the entire compute budget.
  \end{itemize}
\item
  \textbf{Serial-\(N\) multistart}

  \begin{itemize}
  \tightlist
  \item
    the same branches and reducer as parallel-\(N\);
  \item
    executed one after another.
  \item
    Isolates concurrency from breadth.
  \end{itemize}
\item
  \textbf{Parallel-\(N\) best-of-\(N\)}

  \begin{itemize}
  \tightlist
  \item
    independent complete trajectories;
  \item
    no intermediate pruning.
  \item
    Measures simple racing/ensemble benefits.
  \end{itemize}
\item
  \textbf{Parallel adaptive tree}

  \begin{itemize}
  \tightlist
  \item
    short pilots;
  \item
    learned or heuristic evaluation;
  \item
    successive halving;
  \item
    continuation, pruning, and reforking.
  \item
    Tests whether prefix prediction improves over full best-of-\(N\).
  \end{itemize}
\end{enumerate}

All arms should have the same aggregate token/tool/accelerator budget in the primary comparison. The parallel arms alone receive \(N\) reserved lanes. Run a secondary experiment with fixed per-branch budgets to expose the purchased-compute frontier.

\subsection{Suggested primary parameters}\label{suggested-primary-parameters}

\begin{itemize}
\tightlist
\item
  \(N=4\);
\item
  horizon \(X\): one edit/test macrocycle, bounded by tool calls, tokens, and two minutes;
\item
  successive rungs \(X,2X,4X\);
\item
  small exploration reserve;
\item
  fixed total compute budget \(B\);
\item
  fixed deadline \(\tau\).
\end{itemize}

Later ablate:

\[
N\in\{1,2,4,8\},\qquad X\in\{1,3,6\}\text{ macro-actions}
\]

\subsection{Pseudocode}\label{pseudocode}

\begin{Shaded}
\begin{Highlighting}[]
\KeywordTok{def}\NormalTok{ search(task, max\_workers, total\_budget):}
\NormalTok{    root }\OperatorTok{=}\NormalTok{ checkpoint(isolated\_state(task))}
\NormalTok{    frontier }\OperatorTok{=}\NormalTok{ PriorityQueue([root])}
\NormalTok{    active }\OperatorTok{=}\NormalTok{ \{\}}
\NormalTok{    incumbent }\OperatorTok{=} \VariableTok{None}

    \ControlFlowTok{while}\NormalTok{ budget\_remains(total\_budget):}
        \ControlFlowTok{while}\NormalTok{ free\_slots(active, max\_workers) }\KeywordTok{and}\NormalTok{ frontier:}
\NormalTok{            parent }\OperatorTok{=}\NormalTok{ frontier.pop()}

\NormalTok{            proposals }\OperatorTok{=}\NormalTok{ propose\_structured\_options(parent, k}\OperatorTok{=}\DecValTok{4} \OperatorTok{*}\NormalTok{ max\_workers)}
\NormalTok{            proposals }\OperatorTok{=}\NormalTok{ cluster\_and\_select\_diverse\_portfolio(proposals)}

            \ControlFlowTok{for}\NormalTok{ proposal }\KeywordTok{in}\NormalTok{ proposals[:free\_slots(active, max\_workers)]:}
\NormalTok{                child }\OperatorTok{=}\NormalTok{ copy\_on\_write\_fork(parent, proposal)}
\NormalTok{                active[child.}\BuiltInTok{id}\NormalTok{] }\OperatorTok{=}\NormalTok{ launch\_rollout(}
\NormalTok{                    child,}
\NormalTok{                    adaptive\_horizon(child),}
\NormalTok{                    transactional\_tools}\OperatorTok{=}\VariableTok{True}
\NormalTok{                )}

\NormalTok{        result }\OperatorTok{=}\NormalTok{ await\_first\_completion(active)}
\NormalTok{        node }\OperatorTok{=}\NormalTok{ atomic\_checkpoint(result.state)}

\NormalTok{        hard\_signal }\OperatorTok{=}\NormalTok{ run\_public\_tests\_and\_safety\_gates(node)}
\NormalTok{        value, uncertainty, remaining\_cost }\OperatorTok{=}\NormalTok{ evaluate\_prefix(node)}

\NormalTok{        log\_all\_state\_cost\_and\_timing(node)}

        \ControlFlowTok{if}\NormalTok{ public\_acceptance(node):}
\NormalTok{            incumbent }\OperatorTok{=}\NormalTok{ choose\_better(incumbent, node)}

        \ControlFlowTok{if}\NormalTok{ hard\_failure(node):}
\NormalTok{            prune(node)}
        \ControlFlowTok{elif}\NormalTok{ upper\_confidence(value, uncertainty) }\OperatorTok{\textless{}}\NormalTok{ lower\_bound(incumbent):}
\NormalTok{            prune(node)}
        \ControlFlowTok{elif}\NormalTok{ value\_of\_next\_quantum(node) }\OperatorTok{\textgreater{}} \DecValTok{0}\NormalTok{:}
\NormalTok{            frontier.push(node)}
        \ControlFlowTok{elif}\NormalTok{ value\_of\_forking(node) }\OperatorTok{\textgreater{}} \DecValTok{0}\NormalTok{:}
\NormalTok{            frontier.extend(make\_diverse\_children(node))}

\NormalTok{        cancel\_dominated\_active\_branches(active, incumbent)}

\NormalTok{    result }\OperatorTok{=}\NormalTok{ select\_or\_synthesize(frontier, incumbent)}
    \ControlFlowTok{return}\NormalTok{ replay\_in\_clean\_environment\_and\_verify(result)}
\end{Highlighting}
\end{Shaded}

The hidden benchmark grader should not be exposed to the coordinator. Grade every timestamped checkpoint afterward to distinguish:

\begin{itemize}
\tightlist
\item
  \textbf{oracle time-to-first-generated-good branch};
\item
  \textbf{time-to-official-returned-good result};
\item
  selection regret;
\item
  merge-induced regressions.
\end{itemize}

\subsection{Primary metrics}\label{primary-metrics}

\begin{enumerate}
\def\labelenumi{\arabic{enumi}.}
\tightlist
\item
  Official time-to-verified-solution.
\item
  Success by deadline:
  \[
  P(T\le\tau)
  \]
\item
  Restricted mean time-to-solution, including censored failures.
\item
  Total and peak:

  \begin{itemize}
  \tightlist
  \item
    tokens;
  \item
    accelerator-seconds;
  \item
    dollars;
  \item
    tool calls;
  \item
    test CPU time.
  \end{itemize}
\item
  Cost per solved task.
\item
  p50 and p90 latency.
\item
  Coordinator, evaluator, snapshot, and reducer overhead.
\item
  Oracle coverage versus returned-result success.
\item
  Evaluator:

  \begin{itemize}
  \tightlist
  \item
    calibration/Brier score;
  \item
    winner recall;
  \item
    false-prune rate;
  \item
    selected-versus-oracle regret.
  \end{itemize}
\item
  Diversity:
\end{enumerate}

\begin{itemize}
\tightlist
\item
  action and tool trace similarity;
\item
  branch outcome correlation;
\item
  marginal success contributed by each branch.
\end{itemize}

\begin{enumerate}
\def\labelenumi{\arabic{enumi}.}
\setcounter{enumi}{10}
\tightlist
\item
  Canceled-but-billed work and straggler cost.
\end{enumerate}

Use task-paired comparisons, survival curves, and task-cluster bootstrapping. Do not compare latency only among successful runs.

A positive result should require:

\begin{itemize}
\tightlist
\item
  lower time-to-solution or restricted mean latency with confidence intervals;
\item
  non-inferior or improved final success;
\item
  a preregistered compute bound, such as \(1\times\) or \(2\times\) sequential cost;
\item
  advantage over both strong sequential and serial multistart baselines.
\end{itemize}

\subsection{Main experimental failure modes}\label{main-experimental-failure-modes}

\begin{itemize}
\tightlist
\item
  weak or artificially constrained sequential baseline;
\item
  model/API rate limits eliminating real concurrency;
\item
  correlated branches masquerading as \(N\) trials;
\item
  hidden-test leakage;
\item
  flaky benchmark environments;
\item
  coordinator selecting the wrong branch;
\item
  reducer/synthesis destroying a passing result;
\item
  overhead omitted from latency;
\item
  repeated prompts and canceled calls omitted from cost;
\item
  benchmark contamination;
\item
  evaluator overfitting to visible tests;
\item
  interpreting oracle branch coverage as delivered success.
\end{itemize}

\begin{center}\rule{0.5\linewidth}{0.5pt}\end{center}

\section{Conclusion}\label{conclusion}

The future execution of an agent can be made \textbf{partly searchable}, but not perfectly predictable. The most useful abstraction is a portfolio of checkpointed, temporally extended policies operating in isolated environments.

The likely winning design is:

\begin{itemize}
\tightlist
\item
  small, adaptive \(N\), not unlimited replication;
\item
  structured, high-level diversity;
\item
  event-based horizons;
\item
  asynchronous scheduling;
\item
  objective verification;
\item
  uncertainty-aware pruning;
\item
  typed merging;
\item
  transactional side-effect handling;
\item
  explicit latency--compute--quality Pareto accounting.
\end{itemize}

The idea should work best as \textbf{speculative search for uncertain, verifiable tasks}, especially coding, formal reasoning, planning in simulators, and read-only tool research. It is not a general method for shortening inherently sequential cognition, and without a reliable verifier, the coordinator---not generation---will usually become the limiting problem.

\end{document}
